\documentclass[11pt,a4paper]{article}

\usepackage[utf8]{inputenc}
\usepackage[T1]{fontenc}
\usepackage{lmodern}
\usepackage[margin=1in]{geometry}
\usepackage{amsmath,amssymb}
\usepackage{graphicx}
\usepackage{booktabs}
\usepackage{listings}
\usepackage{xcolor}
\usepackage{hyperref}
\usepackage{url}
\usepackage{natbib}
\usepackage{authblk}

\hypersetup{
    colorlinks=true,
    linkcolor=blue!60!black,
    citecolor=blue!60!black,
    urlcolor=blue!60!black,
}

\lstset{
    basicstyle=\ttfamily\footnotesize,
    breaklines=true,
    frame=single,
    backgroundcolor=\color{gray!5},
    showstringspaces=false,
}

\title{\texttt{pixelflow}: GPU Texture Reservoir Computing via Fragment-Shader Cellular Automata}

\author{Francisco Angulo de Lafuente}
\affil{Independent Researcher \\ \texttt{lareliquia.angulo@gmail.com}}

\date{\today}

\begin{document}
\maketitle

\begin{abstract}
We present \texttt{pixelflow}, an open-source Python library for reservoir
computing in which a 2D GPU texture serves as the reservoir substrate. Each
pixel holds a four-channel state and evolves under a local cellular-automaton
rule expressed as a fragment shader; the flattened final texture is used as a
feature vector for a linear readout. This is a software analogue of physical
reservoir computing and exploits the GPU's native strength in massively
parallel per-pixel updates with local neighbourhood access. The library
provides a CPU reference backend (NumPy) and a headless GPU backend (moderngl,
OpenGL~3.3 core) that agree numerically within $10^{-5}$ absolute error. Three
built-in cellular-automaton rules (Gray-Scott reaction-diffusion, continuous
Life, and a discrete wave equation) and three input encoders (tile, phase,
random projection) enable experimentation. We report honest benchmarks on
MNIST digit classification and a synthetic two-moons task, including cases
where the reservoir under-performs a raw-pixel linear baseline. \texttt{pixelflow}
is released under Apache-2.0 at
\url{https://github.com/Agnuxo1/pixelflow}.
\end{abstract}

\section{Introduction}

Reservoir computing (RC) separates a dynamical system from a trainable
readout: inputs are injected into a high-dimensional, generally random,
dynamical reservoir, and only a simple linear model is fit to the reservoir's
state \citep{jaeger2001,maass2002}. This paradigm has been instantiated in
software (echo state networks, liquid state machines) and in physical
substrates, including optical systems, memristive arrays, and spintronic
devices \citep{tanaka2019}. Physical RC is attractive because it offloads the
computationally expensive part (the recurrent dynamics) onto native hardware.

This work proposes a GPU texture as a software-level reservoir. Modern GPUs
are optimised for exactly the operations a cellular-automaton (CA) reservoir
requires: sampling a local neighbourhood on a 2D grid and writing a per-pixel
update. A fragment shader implementing a local rule expresses this directly,
and ping-pong framebuffers provide efficient temporal iteration. The idea is
simple, and the novelty here is not in the mathematics of CA reservoirs but
in a clean, tested, and honest reference implementation that can serve as a
baseline for further research.

\section{Method}

\subsection{State and dynamics}

Let $S_t \in \mathbb{R}^{H \times W \times 4}$ denote the reservoir state at
discrete time $t$. The initial state $S_0$ is derived from an input
$\mathbf{x} \in \mathbb{R}^d$ by one of three encoders:

\begin{itemize}
    \item \textbf{tile}: $\mathbf{x}$ is reshaped to a $(H, W)$ grid
    (nearest-neighbour resize) and placed in channel~0; remaining channels are
    filled with small random noise.
    \item \textbf{phase}: each scalar $x_i$ is mapped to $(\sin x_i, \cos x_i)$
    in channels~0 and~1.
    \item \textbf{project}: a fixed random Gaussian matrix
    $W \in \mathbb{R}^{(HWC) \times d}$, seeded for reproducibility, projects
    $\mathbf{x}$ to the full state tensor.
\end{itemize}

The update $S_{t+1} = \Phi(S_t)$ is implemented as a local operator with
periodic boundary conditions. Three rule families are provided:

\paragraph{Gray-Scott reaction-diffusion} Two channels encode concentrations
$u, v$ obeying
\begin{align}
    \frac{\partial u}{\partial t} &= D_u \nabla^2 u - u v^2 + F(1 - u), \\
    \frac{\partial v}{\partial t} &= D_v \nabla^2 v + u v^2 - (F + k) v,
\end{align}
discretised with a 5-point Laplacian and default parameters
$F=0.055$, $k=0.062$, $D_u=0.16$, $D_v=0.08$, $\mathrm{d}t=1.0$.

\paragraph{Continuous Life} A smoothed Conway's Life rule applied to channel
0: a local average over a 3$\times$3 neighbourhood is passed through a
threshold $\tau$ with optional additive noise.

\paragraph{Discrete wave equation} Channel 0 holds amplitude $u$ and channel
1 holds velocity $\dot u$, updated as
\begin{align}
    \dot u_{t+1} &= \gamma \dot u_t + c^2 \nabla^2 u_t, \\
    u_{t+1} &= u_t + \mathrm{d}t \cdot \dot u_{t+1},
\end{align}
with damping $\gamma = 0.999$ by default.

\subsection{Readout}

After $k$ updates the feature vector $\phi(\mathbf{x}) = \mathrm{vec}(S_k) \in
\mathbb{R}^{HWC}$ is fed to a linear model from scikit-learn
\citep{pedregosa2011}: \texttt{Ridge}, \texttt{RidgeClassifier}, or
\texttt{LogisticRegression}. Training is closed-form (Ridge) or convex
(Logistic); no gradient flows through the reservoir.

\subsection{Backends}

Two interchangeable backends implement $\Phi$:

\begin{itemize}
    \item \textbf{CPU}: vectorised NumPy \citep{harris2020} using
    \texttt{np.roll} for periodic neighbourhood access.
    \item \textbf{moderngl}: headless OpenGL~3.3 core context,
    \texttt{GL\_RGBA32F} textures, ping-pong framebuffers, one fragment shader
    per rule.
\end{itemize}

Parity is verified by \texttt{tests/test\_moderngl\_backend.py}: for each rule,
the CPU and GPU outputs on an $8{\times}8{\times}4$ random state after 5 steps
agree to $< 10^{-5}$ absolute error. GPU determinism across repeated runs is
also tested.

\section{Experiments}

All numbers below are produced by scripts in the repository's
\texttt{benchmarks/} directory and stored as JSON with metadata (hardware,
config, timings) in \texttt{benchmarks/results/}. We report raw numbers,
including cases where the reservoir is outperformed by a raw-pixel linear
baseline.

\subsection{Synthetic: two-moons}

A 500-training / 200-test two-moons dataset ($\mathrm{noise}=0.2$) with a
$16{\times}16{\times}4$ reservoir, 5 steps of \texttt{diffusion\_reaction} and
\texttt{project} encoding, followed by \texttt{RidgeReadout}, reaches
test accuracy \textbf{0.92}.

\subsection{MNIST (subset 3\,000)}

Quick smoke runs with a $32{\times}32{\times}4$ reservoir:

\begin{table}[h]
\centering
\begin{tabular}{lccc}
\toprule
Rule           & Steps & Test acc. & Baseline (raw+LogReg) \\
\midrule
\texttt{wave}              & 4 & 0.879 & 0.886 \\
\texttt{life\_like}        & 6 & 0.348 & 0.886 \\
\bottomrule
\end{tabular}
\caption{MNIST 3\,000-sample subset, tile encoding. The \texttt{life\_like}
rule saturates and destroys the digit signal at these parameters; the
\texttt{wave} rule preserves it and nearly matches the baseline.}
\end{table}

\subsection{MNIST (full 60\,000 / 10\,000 split)}

\begin{table}[h]
\centering
\begin{tabular}{lcccc}
\toprule
Rule                  & Grid          & Steps & Test acc. & Baseline (raw+LogReg) \\
\midrule
\texttt{wave}              & $32{\times}32{\times}4$ & 4 & \textbf{0.9281} & 0.9261 \\
\texttt{diffusion\_reaction} & $32{\times}32{\times}4$ & 8 & 0.9213 & 0.9261 \\
\bottomrule
\end{tabular}
\caption{MNIST full 60k training / 10k test split, tile encoding, \texttt{cpu}
backend, \texttt{RidgeClassifier} readout. Raw JSON artifacts with timings and
hardware metadata are in \texttt{benchmarks/results/}. The \texttt{wave}
reservoir slightly beats the raw-pixel baseline; \texttt{diffusion\_reaction} at
the chosen parameters slightly underperforms it. No hyperparameter search was
run: these are first-order honest numbers.}
\label{tab:mnist60k}
\end{table}

\subsection{Eikonal equation}

The repository includes a reference fast-marching Eikonal solver against which
the \texttt{wave} reservoir's arrival-time field can be compared. A full
quantitative study is left for future work.

\section{Limitations}

\begin{enumerate}
    \item \textbf{No attempt to beat deep CNNs.} A standard CNN exceeds 99\%
    accuracy on MNIST; the point of \texttt{pixelflow} is dynamics and
    reproducibility, not leaderboard chasing.
    \item \textbf{GPU overhead dominates for small inputs.} The moderngl
    backend is slower than CPU for batches $< 1000$ on $< 32{\times}32$
    grids due to standalone-context creation. A shared-context v0.2 mitigates
    this; the current version accepts the honest cost.
    \item \textbf{Fixed CA rules.} Learned reservoirs, in the style of
    \citet{mordvintsev2020}, are future work.
\end{enumerate}

\section{Reproducibility}

Every benchmark script prints its config and writes a timestamped JSON with
the exact parameters, random seed, and hardware identification. Seeds are
fixed in all tests and examples. The full test suite
(\texttt{pytest tests/}) comprises 44~tests that run on CI across
Linux/macOS/Windows and Python~3.10/3.11/3.12.

\section{Conclusion}

\texttt{pixelflow} distils a single clean idea --- fragment-shader CA as a
reservoir substrate with a linear readout --- into a tested, honest, and
reusable library. It is not a new deep-learning state of the art, and it is
not claimed to be. It is a clean testbed for reservoir-computing research
that happens to run fast on commodity GPUs.

\section*{Availability}

Source, tests, benchmarks, and documentation are available at
\url{https://github.com/Agnuxo1/pixelflow} under the Apache-2.0 licence.

\begin{thebibliography}{10}

\bibitem[Jaeger(2001)]{jaeger2001}
Jaeger, H.\ (2001). \textit{The ``echo state'' approach to analysing and
training recurrent neural networks.} GMD Report 148, German National Research
Center for Information Technology.

\bibitem[Maass et~al.(2002)]{maass2002}
Maass, W., Natschläger, T., \& Markram, H.\ (2002). Real-time computing without
stable states: A new framework for neural computation based on perturbations.
\textit{Neural Computation}, 14(11), 2531--2560.

\bibitem[Tanaka et~al.(2019)]{tanaka2019}
Tanaka, G., Yamane, T., Héroux, J.~B., Nakane, R., Kanazawa, N., Takeda, S.,
Numata, H., Nakano, D., \& Hirose, A.\ (2019). Recent advances in physical
reservoir computing: A review. \textit{Neural Networks}, 115, 100--123.

\bibitem[Pedregosa et~al.(2011)]{pedregosa2011}
Pedregosa, F.\ et al.\ (2011). Scikit-learn: Machine learning in Python.
\textit{JMLR}, 12, 2825--2830.

\bibitem[Harris et~al.(2020)]{harris2020}
Harris, C.~R.\ et al.\ (2020). Array programming with NumPy.
\textit{Nature}, 585, 357--362.

\bibitem[Mordvintsev et~al.(2020)]{mordvintsev2020}
Mordvintsev, A., Randazzo, E., Niklasson, E., \& Levin, M.\ (2020). Growing
neural cellular automata. \textit{Distill}.

\end{thebibliography}

\end{document}
